\guidechapter{What a successful check actually tells us}{Three claims that should not be confused.}

There are different meanings of ``it works.'' A circuit may return the
right values on a specified finite set of inputs. A physical program may
execute the circuit with a proved resource bound. A theorem may cover every
input length, including sizes no experiment can run. Each is valuable, but
they are different statements.

\explfigure{23-proof-obligations}{Read the claim at the level its evidence supports.}{
The columns separate finite algebra, physical implementation, and all-size
analysis. They are obligations to connect, not competing ways of proving the
same statement.}

\paragraph{A concrete example.}
The register plot shows that the four updates work when the scratch value
starts at seven. The symbolic calculation proves restoration for arbitrary
values in its stated ring. Neither statement tells us how much tape
movement an enormous array implementation needs. That requires another
argument. Similarly, verifying the seven rational margins checks arithmetic
within a cost model; it does not establish every algorithmic premise of that
model.

\paragraph{Why the guide labels toy examples.}
Large, visible exponent gaps are useful for learning but would misrepresent
the paper if presented as its measured speedup. Small matrices and arrays
can demonstrate an invariant without being the actual finite construction.
The captions say which numbers come from the source and which are teaching
examples. The new figure generators check the displayed arithmetic so that
the illustrations explain real identities rather than suggestive pictures.

\paragraph{After this guide.}
You should be able to describe why convolution helps, why moving data costs
time, why common frames preserve addition, how a weighted recursion saves
an exponent, and why recovery requires a strict error budget. Following
every detail of the source theorem still requires further study of linear
algebra, Fourier analysis, and complexity proofs. The expanded edition puts
the original arguments immediately after their explanatory material.

\takeaway{Understanding the architecture is a substantial step. It is
not the same claim as having independently verified every part of the theorem.}
